A large class of production decisions is not open-ended generation but a
fixed menu: route this message to one of 77 intents, flag whether it
mentions clinical risk, rate its urgency on a four-point scale. Language
models answer such questions by generating text that must then be parsed
and constrained, at a cost per question and with no guarantee the answer
lies on the menu.

We describe \openjev{}, an open architecture for \emph{typed decisions}. A
state is encoded once as a shared prefix; every candidate option is placed
after it behind a marker token; block-diagonal attention isolates the
questions from one another; and a grouped log-softmax normalises scores
within each question's own option set. Ten questions about one input
therefore cost one forward pass, and the answer is a calibrated probability
over a closed menu rather than generated text. Latency is flat in the
number of questions: ten cost 6\% more than one.

Our central finding is that this architecture, on its own, does not
generalize. A strong bidirectional encoder with the output format and no
task training scores \textbf{0.7$\times$ chance} on held-out schemas, and a
61-task training mixture leaves it at 0.9$\times$. Generalization comes from
the composition of the training data, and specifically from properties that
are cheap to synthesize: padding training menus with borrowed distractor
labels takes accuracy on a 77-way held-out menu from a literal $0.000$ to
$0.205$, because a model that has never seen a large menu collapses onto one
label when handed one. Scaling an audited mixture to 279 tasks and 323{,}466
rows brings the held-out suite to \textbf{21.9$\times$ chance}, with all
seven tasks clearing their chance floor and a 151-way menu reaching
94.9$\times$.

Given labels for the target task, small models are competitive with a
commercial typed-decision API. A rank-16 LoRA adapter on Qwen3-1.7B, trained
on 395 examples in 258 seconds, \textbf{beats} that API on a healthcare
routing benchmark: intent $0.979$ against $0.941$ ($p=1\mathrm{e}{-3}$),
multi-label exact set $0.909$ against $0.822$ ($p=7\mathrm{e}{-6}$), and
three simultaneous intents $0.903$ against $0.645$; three wins, four
statistical ties, no losses. Notably the adapter, at $17$M trainable
parameters and $87$\,MB, outperforms opening all $1{,}725$M ($0.929$,
$3.4$\,GB).

We also report where this fails. Without labels the gap remains large:
$0.290$ on Banking77 against roughly $0.820$. And benchmark transfer is not
deployment readiness. The same general checkpoint scoring 21.9$\times$
chance, applied zero-shot to the routing task, reaches intent $0.601$ and a
clinical-gate recall of $0.111$. Every claim here is reproducible from a
public repository, with datasets, checkpoints, and the negative results
reported alongside the positive ones.
